\documentclass[11pt]{article}
\usepackage[T1]{fontenc}
\usepackage[utf8]{inputenc}
\usepackage[margin=1.05in]{geometry}
\usepackage{amsmath,amssymb,amsthm,mathtools}
\usepackage{lmodern,microtype}
\usepackage[authoryear,round]{natbib}
\usepackage{xurl}
\usepackage[colorlinks=true,allcolors=blue]{hyperref}
\newcommand{\workbenchurl}{https://github.com/KokunoYumeto/schemes-and-moduli-workbench}
\newcommand{\recordurl}{https://github.com/KokunoYumeto/mathematics-workbenches}
\newcommand{\OO}{\mathcal O}\newcommand{\Fp}{\mathbf F_p}
\newcommand{\etale}{\'{e}tale}
\DeclareMathOperator{\Spec}{Spec}\DeclareMathOperator{\Hom}{Hom}
\DeclareMathOperator{\QCoh}{QCoh}\DeclareMathOperator{\Crys}{Crys}
\DeclareMathOperator{\Strat}{Strat}
\newtheorem{theorem}{Theorem}[section]
\newtheorem{lemma}[theorem]{Lemma}\newtheorem{proposition}[theorem]{Proposition}
\newtheorem{corollary}[theorem]{Corollary}
\theoremstyle{definition}\newtheorem{definition}[theorem]{Definition}
\theoremstyle{remark}\newtheorem{remark}[theorem]{Remark}
\title{Grothendieck's comparison of infinitesimal crystals and ordinary stratifications}
\author{}\date{}
\begin{document}\maketitle
\begin{abstract}
Grothendieck proposed extending the comparison between infinitesimal crystals and stratified modules from smooth morphisms to flat morphisms of finite presentation, with a proviso on Noetherian bases and coherent coefficients. We prove this qualified statement in every characteristic, using Jiang's derived comparison and a separate passage to ordinary coherent modules. We also prove the comparison for arbitrary flat coefficients over a Noetherian base. Two explicit counterexamples explain why the qualifications matter: over a non-Noetherian base, even an automorphism of the structure crystal can be invisible on the underlying scheme; over a regular Noetherian base the same phenomenon occurs with a noncoherent target. In characteristic $p$, we describe arbitrary crystals by relative Frobenius-divided modules and obtain positive results for pure Frobenius and cuspidal normalizations. The mathematics and text were produced by GPT-6 Astra Pro in ChatGPT Pro mode and checked by targeted reading and a separate statement-only re-proof of all six main results by the same model; no human has checked them.
\end{abstract}

\section{The question and the comparison theorems}

An infinitesimal crystal specifies a module on every nilpotent thickening and makes those modules compatible with pullback. A stratification records parallel transport only between infinitesimally close points of the scheme itself. Recovering the first from the second is immediate for smooth schemes, because maps lift locally across nilpotent immersions. For a singular flat scheme such lifts can fail, so this recovery is a substantive descent question.

In \S1.8 of the notes collected as cote~6 of the Montpellier archive, Grothendieck asks for the comparison when the morphism is flat and locally of finite presentation. His qualification reads ``du moins si R loc.noeth. et se bornant aux Modules coh\'erents'' \citep[archive sheet~6, PDF image~7]{GrothendieckNotes}. We use $S$ for his base and $X$ for his scheme. Theorem~\ref{thm:coherent} proves the qualified assertion. Theorems~\ref{thm:counterexample} and Corollary~\ref{cor:noetherian-counterexample} explain what happens when either finiteness requirement is removed. They concern the broader assertion, not a contradiction to the qualified sentence.

\begin{definition}
On the ordinary infinitesimal site $(X/S)_{\mathrm{inf}}$, an object is a nilpotent closed immersion $U\hookrightarrow T$ over $S$, where $U$ is open in $X$. A quasi-coherent crystal $\mathcal E$ consists of quasi-coherent $\OO_T$-modules $\mathcal E_T$ and compatible isomorphisms $h^*\mathcal E_{T'}\simeq\mathcal E_T$ for every morphism of tests. All pullbacks in this definition are ordinary tensor pullbacks. In the big site, $U$ is allowed to be an arbitrary $X$-scheme. The topology is Zariski.

An ordinary stratification on $M\in\QCoh(X)$ is a compatible isomorphism between its two pullbacks to every finite infinitesimal neighborhood of the diagonal in $X\times_SX$. It is the identity on the diagonal and satisfies the cocycle on the triple neighborhoods. For a nonseparated morphism these neighborhoods are formed locally around the diagonal. These local definitions glue. The evaluation of a crystal on $X$, together with its diagonal transition maps, defines a functor
\[
 \mathrm{res}:\Crys(X/S)\longrightarrow\Strat(X/S).
\]
Neither definition includes divided powers.
\end{definition}

\begin{theorem}[Ordinary coherent crystals for flat morphisms over a Noetherian base]\label{thm:coherent}
Let $S$ be locally Noetherian and let $X\to S$ be flat and locally of finite presentation.  Restriction is an equivalence
\[
 \left\{\begin{array}{c}
 \text{quasi-coherent infinitesimal crystals}\\
 \text{with coherent value on }X
 \end{array}\right\}
 \ \simeq\
 \left\{\begin{array}{c}
 \text{ordinary coherent stratified}\\
 \OO_X\text{-modules over }S
 \end{array}\right\}.
\]
The assertion uses ordinary tensor pullback in both categories.  It holds for the small or big infinitesimal site, also when arbitrary classical tests are allowed.  On Noetherian tests the values of these crystals are coherent; on arbitrary tests they are of finite presentation.  No flatness of the coefficient module, and no restriction on the characteristic, is imposed.


\end{theorem}

\begin{theorem}\label{thm:flat}
Let $S$ be locally Noetherian and let $X\to S$ be flat and locally of finite presentation.  Restriction is an equivalence between quasi-coherent crystals whose values on infinitesimal tests are flat modules and ordinary stratified quasi-coherent $\OO_X$-modules that are flat over $\OO_X$.  There is no finite-generation hypothesis on the coefficients, and no restriction on the characteristic of $S$.

This equivalence restricts to finite locally free coefficients.  It holds for either the small or big infinitesimal site, with arbitrary classical test schemes.
\end{theorem}

\begin{theorem}\label{thm:counterexample}
For every prime $p$ there is an $\Fp$-algebra $R$ and a finite free morphism $X\to\Spec R$ of rank $p$ for which the restriction functor
\[
 \{\text{crystals of $\OO$-modules on $(X/R)_{\rm inf}$}\}
 \longrightarrow
 \{\text{ordinary stratified $\OO_X$-modules relative to $R$}\}
\]
is not faithful.  The failure already occurs among automorphisms of the rank-one locally free structure crystal.  In particular, flatness and finite presentation of the morphism do not by themselves make restriction faithful for all coefficients.
\end{theorem}

The coefficients in Theorem~\ref{thm:coherent} may be nonflat. For example, a coherent module on the base defines a constant-base crystal even when it has torsion. Theorem~\ref{thm:flat} addresses a different direction: its modules can have infinite rank. In Theorem~\ref{thm:counterexample}, $X/S$ is finite free and every fiber is a complete intersection. Thus the obstruction does not come from a failure of flatness or from a badly behaved geometric fiber.

In positive characteristic there is a second description that works without either finiteness qualification on the base or the coefficient module. A ring map $A\to B$ is \emph{pure} if $N\to B\otimes_A N$ is injective for every $A$-module $N$; purity here is a module-theoretic condition.

\begin{theorem}\label{thm:frobenius}
Let $R$ be any scheme of characteristic $p$, and let $X\to R$ be locally of finite presentation.  Put $X^{(n)}=X\times_{R,F_R^n}R$.  Quasi-coherent crystals on $(X/R)_{\rm inf}$ are equivalent to systems of arbitrary quasi-coherent modules
\[
 M_n\in\operatorname{QCoh}(X^{(n)}),\qquad
 F_{X^{(n)}/R}^*M_{n+1}\simeq M_n.
\]
If every relative Frobenius $F_{X^{(n)}/R}$ is pure, restriction gives an equivalence between these crystals and ordinary stratified quasi-coherent $\OO_X$-modules.  No perfectness, reducedness, or Noetherian hypothesis on $R$ is required.
\end{theorem}

\subsection{Published inputs and the point of the argument}

Grothendieck's published account distinguishes the infinitesimal and stratifying sites and gives their agreement under formal smoothness \citep[\S4.1--4.2, pp.~331--333]{Grothendieck1968}. The distinction is essential in the counterexamples below. For singular schemes, a comparison of derived categories need not identify their ordinary coefficient categories without further argument.

Jiang proves a derived comparison for a separated finite-type morphism of finite Tor-amplitude over a Noetherian base \citep[Theorem~2.1, Remark~2.4, Proposition~6.2 and Theorem~6.3, pp.~16--17 and 32--34]{Jiang}. We use that theorem, not a claim that all its derived objects are ordinary modules. Flatness of $X/S$ makes the completed diagonal rings flat over every chosen factor. Coherent stratifications therefore lift to the derived formal nerve. On a nilpotent test the resulting derived crystal is concentrated in nonpositive cohomological degrees; taking $H^0$ then commutes with ordinary pullback. A completed ordinary \v{C}ech calculation separately proves full faithfulness.

For flat coefficients a finite filtration by powers of the nilpotent ideal shows more: every derived test value is an ordinary flat module. This is why infinite rank causes no difficulty in Theorem~\ref{thm:flat}. The counterexample uses a module with principal parts of unbounded order. Its compatible Frobenius-divided sections vanish on $X$ and survive on a thicker test, precisely where finite generation and the Noetherian separation argument cease to apply.

The Frobenius description is proved directly by descent from affine space. Purity then makes all the resulting module descent data effective, using the theorem of Joyal--Tierney published by \citet[Theorem~4, p.~41]{Mesablishvili}. The cusp conclusion agrees with the ordinary-pullback equivalence of \citet[Proposition~2.11, Theorems~3.11 and~4.3, Remark~4.5]{BenZviNevins}. We give a direct proof using a finite faithfully flat reverse Frobenius factor and ordinary descent; it shows that purity is sufficient but not necessary. The local freeness of coherent coefficients over a field agrees with \citet[Theorem~6.8(3)--(4), arXiv pp.~22--25]{ToniniZhang}. We make no priority claim for consequences of these published inputs. The proofs below establish the stated ordinary coefficient scope and give the explicit failures outside it.

Recent AI-assisted work on the same archive includes Hua's construction of the universal truncated divided-power algebra, with a reported Lean verification \citep[\S35, Theorems~35.1--35.3, pp.~161--162]{Hua}. It concerns the divided-power foundations of crystalline theory. The present comparison concerns ordinary nilpotent tests and ordinary tensor pullbacks; its proofs and verification status are given explicitly below.

\section{Completed descent for ordinary coherent modules}

This section proves that a morphism of coherent crystals is recovered from its stratified value and prepares the ordinary part of the comparison. Completion is taken only over Noetherian rings and on finite modules when tensor products and inverse limits are compared.

\begin{lemma}[An exact completed \v{C}ech complex over a nilpotent test]\label{lem:cech}
Let $R$ be Noetherian, let $A$ be a flat finitely presented $R$-algebra, and let $B$ be a Noetherian $R$-algebra with a nilpotent ideal $I$.  Put $C=B/I$ and fix an $R$-algebra map $a:A\to C$.  For $q\geq0$ put
\[
 D^q=B\otimes_R A^{\otimes_R(q+1)},\qquad
 J_q=\ker(D^q\longrightarrow C),\qquad
 P^q=\varprojlim_n D^q/J_q^n,
\]
where the map to $C$ uses $B\to C$ and $a$ on every $A$-factor.  The usual coface maps, inserting $1$ in an $A$-factor, define an augmented cosimplicial $B$-algebra
\[
                    B\longrightarrow P^0\rightrightarrows P^1
                              \longrightarrow P^2\longrightarrow\cdots.
\]
For every $B$-module $N$ its alternating augmented complex
\[
 0\longrightarrow N\longrightarrow P^0\otimes_BN
       \longrightarrow P^1\otimes_BN\longrightarrow\cdots
\]
is exact.  In particular the first two arrows exhibit $N$ as the equalizer of the two coface maps from $P^0\otimes_BN$ to $P^1\otimes_BN$.
\end{lemma}
\begin{proof}
Every $D^q$ is a Noetherian ring and is flat over $B$.  Noetherian adic completion is flat, so every $P^q$ is flat over $B$.  Completion is exact on finite modules and commutes with their tensor description; consequently
\[
 P^q/IP^q\simeq
 \widehat{C\otimes_R A^{\otimes_R(q+1)}},
\]
where the completion on the right is along the graph ideal.  These are the precise Noetherian completion facts in \citet[Lemmas~10.97.1--6, especially Tags~00MA and~00MB]{StacksProject}.

Modulo $I$, the augmented formal \v{C}ech nerve has a section.  Explicitly, applying $a$ to the first $A$-factor and multiplying its image into the $C$-factor gives continuous homomorphisms
\[
 C\otimes_R A^{\otimes_R(q+1)}
       \longrightarrow C\otimes_R A^{\otimes_R q}.
\]
They preserve the graph ideals and hence extend to the completions.  For $q=0$ the target is $C$.  These maps are extra codegeneracies: their identities with insertion of $1$ are the usual split augmented \v{C}ech identities.  The resulting alternating complex is therefore contractible after augmentation, and remains so after tensoring over $C$ with any $C$-module.

Filter $N$ by the finite filtration $I^jN$.  Since all $P^q$ are $B$-flat, this induces a filtration by subcomplexes whose successive quotient for $Q_j=I^jN/I^{j+1}N$ is the just-described split complex
\[
 0\longrightarrow Q_j\longrightarrow
 (P^0/IP^0)\otimes_CQ_j\longrightarrow
 (P^1/IP^1)\otimes_CQ_j\longrightarrow\cdots.
\]
Every quotient is exact.  Induction on the finite filtration, using the long exact cohomology sequence, proves exactness of the original complex.  No finiteness condition on $N$ was used.
\end{proof}

\begin{proposition}\label{prop:faithfulness}
Let $S$ be locally Noetherian and let $X\to S$ be flat and locally of finite presentation.  Let $\mathcal E,\mathcal F$ be quasi-coherent crystals on the ordinary infinitesimal site of $X/S$, with $\mathcal F_X$ coherent.  Then evaluation on $X$ is injective on morphisms:
\[
 \Hom_{(X/S)_{\rm inf}}(\mathcal E,\mathcal F)
       \longrightarrow \Hom_{\OO_X}(\mathcal E_X,\mathcal F_X).
\]
In particular the restriction functor from coherent crystals to ordinary coherent stratified modules is faithful.  The assertion holds for the small and the big infinitesimal site, including when arbitrary, possibly non-Noetherian, test schemes are allowed.  No finiteness assumption on $\mathcal E$ is required.
\end{proposition}
\begin{proof}
Let $u:\mathcal E\to\mathcal F$ have $u_X=0$.  We first work over an affine base $S=\Spec R$ and an affine open $V=\Spec A\subset X$ with $A$ a finitely presented flat $R$-algebra.  Consider an affine Noetherian test $U=\Spec C\hookrightarrow T=\Spec B$, whose structural map $U\to X$ factors through $V$.  Thus $B\twoheadrightarrow C$ has a nilpotent kernel $I$, and there is an $R$-algebra map $A\to C$.  In the small site, after restricting to affine opens, we may take $U=V$.

The crystal property identifies $\mathcal F_T/I\mathcal F_T$ with the pullback of $\mathcal F_X$ to $U$, which is a finite $C$-module.  Lifting a finite set of generators and iterating $I^N=0$ shows that $\mathcal F_T$ is a finite $B$-module; this elementary nilpotent form of Nakayama's argument does not require $\mathcal F_T$ to be finite in advance.

Set
\[
 D=A\otimes_R B,\qquad
 J=\ker\bigl(D\longrightarrow C,\ a\otimes b\longmapsto a\overline b\bigr).
\]
The map to $C$ is surjective.  The ring $D$ is Noetherian, since it is a finitely presented $B$-algebra, and it is flat over $B$, since $A$ is flat over $R$.  For every $n\geq1$, the closed immersion
\[
 U\hookrightarrow T_n=\Spec(D/J^n)
\]
is an infinitesimal test, with compatible maps to $T$ and to $V$.  The latter map and $u_X=0$ give $u_{T_n}=0$.  The former map and the crystal transition isomorphisms therefore give
\[
 (D/J^n)\otimes_B u_T=0\qquad(n\geq1).
\]
Put $N=D\otimes_B\mathcal F_T$, a finite $D$-module.  Every element in the image of $D\otimes_B u_T$ belongs to $\bigcap_nJ^nN$.

Fix any prime $\mathfrak q\subset B$.  Since $I$ is nilpotent, it determines a prime of $C=B/I$ and hence a prime $\mathfrak p\subset D$ containing $J$, with $\mathfrak p\cap B=\mathfrak q$.  Krull's intersection theorem in the Noetherian local ring $D_{\mathfrak p}$ gives
\[
 \bigcap_{n\geq1}(JD_{\mathfrak p})^nN_{\mathfrak p}=0.
\]
Thus $D_{\mathfrak p}\otimes_Bu_T=0$.  The map $B_{\mathfrak q}\to D_{\mathfrak p}$ is flat and local, hence faithfully flat.  Consequently $(u_T)_{\mathfrak q}=0$.  As $\mathfrak q$ was arbitrary, $u_T=0$.  The precise commutative-algebra inputs are Krull intersection and faithful flatness of a flat local map, \citet[Lemmas~10.51.4 and~10.39.17, Tags~00IP and~00HR]{StacksProject}; the slightly more general separation statement is Lemma~10.51.5, Tag~00IQ.

It remains to remove the Noetherian condition on the test.  Locally write $A=R[z_1,\ldots,z_r]/K$, lift the images of the $z_i$ from $C$ to $B$, and choose $N$ with $I^N=0$.  The resulting map $R[z_1,\ldots,z_r]\to B$ annihilates $K^N$, and hence factors through the Noetherian ring
\[
 B_0=R[z_1,\ldots,z_r]/K^N.
\]
The test $(U,T)$ maps to the finite-type test
$V\hookrightarrow\Spec B_0$.  In the small site $U$ is an open of $V$, as required; in the big site its given map to $V$ suffices.  We already proved $u_{\Spec B_0}=0$.  The crystal transition isomorphisms then give $u_T=0$.  These local conclusions glue over $X$ and $S$, proving the proposition.
\end{proof}

\section{Coherent and flat coefficients}

We now apply the derived comparison in a form that records the difference between derived and ordinary pullback. We write $D(B)$ for the unbounded derived category of $B$-modules, with cohomological grading; $D^{\leq0}(B)$ consists of complexes with no positive cohomology. Derived tensor product is right t-exact in this convention.

\begin{proof}[Proof of Theorem~\ref{thm:coherent}]
We divide the proof into ordinary full faithfulness and essential surjectivity.  Only the latter uses a derived comparison theorem.  In that step we take zeroth cohomology; we do not assert that the intermediate derived crystal is concentrated in degree zero.

\emph{Full faithfulness on Noetherian tests.}
Let $\mathcal E,\mathcal F$ be the crystals in the statement and let
$f_X:\mathcal E_X\to\mathcal F_X$ be a morphism of their ordinary stratified modules.  Work over $S=\Spec R$ and an affine open $V=\Spec A\subset X$.  On an affine Noetherian test $U=\Spec C\hookrightarrow T=\Spec B$ whose map to $X$ factors through $V$, retain the notation $I,D^q,J_q,P^q$ from the lemma.  In the small site we may take $U=V$ after localization.  Since $I$ is nilpotent, lifting a finite set of generators modulo $I$ shows that both $E_T=\Gamma(T,\mathcal E_T)$ and $F_T=\Gamma(T,\mathcal F_T)$ are finite $B$-modules.  Hence they are finitely presented.

For every $q,n$, the test $U\hookrightarrow\Spec(D^q/J_q^n)$ maps both to $T$ and, using any selected $A$-factor, to $V$.  The crystal identifications, followed by completion of finite modules, give canonical isomorphisms
\begin{equation}\label{eq:formal-values}
 P^q\otimes_BE_T\simeq P^q\otimes_AE_V,\qquad
 P^q\otimes_BF_T\simeq P^q\otimes_AF_V.
\end{equation}
Here the chosen $A$-factor specifies the right-hand tensor product.  This passage to completion is justified because $D^q\otimes_BE_T$ and $D^q\otimes_AE_V$, and likewise the two modules for $F$, are finite over the Noetherian ring $D^q$.  No tensor product with an arbitrary inverse limit is interchanged.

Use \eqref{eq:formal-values} for $q=0$ to transport $P^0\otimes_Af_V$ to a homomorphism
\[
                 \widetilde f:P^0\otimes_BE_T\longrightarrow P^0\otimes_BF_T.
\]
Its two pullbacks to $P^1$ agree.  Indeed the two choices of $A$-factor in \eqref{eq:formal-values} differ by the given stratification: this follows already on each finite graph neighbourhood from the crystal cocycle, and then by completion.  Compatibility of $f_X$ with these stratifications is exactly the required equality.

Put $H=\Hom_B(E_T,F_T)$, a finite $B$-module.  Since $E_T$ is finitely presented and $P^q$ is $B$-flat, the canonical map
\[
 P^q\otimes_BH\ \simeq\
 \Hom_{P^q}(P^q\otimes_BE_T,P^q\otimes_BF_T)
\]
is an isomorphism; it follows by applying $\Hom(-,F_T)$ to a finite presentation of $E_T$ and using flatness to preserve its kernel.  The equalizer assertion of Lemma~\ref{lem:cech} therefore descends $\widetilde f$ to a unique $B$-linear map $f_T:E_T\to F_T$.

These maps are compatible with morphisms and open restrictions of Noetherian tests.  To check this, pull both proposed maps to the completion $P^0$ for the source test.  The maps of finite graph neighbourhoods, and hence their completions, identify both pullbacks with the given map $f_V$.  The map $B\to P^0$ is faithfully flat: it is flat as above and its composition with $P^0\twoheadrightarrow C$ is the nilpotent quotient $B\twoheadrightarrow C$, whose spectrum is surjective.  The two maps on the test are consequently equal.  The same equalizer uniqueness shows that a morphism of crystals is uniquely determined by its value on $X$.  This proves ordinary full faithfulness on Noetherian tests.

\emph{Lifting a coherent stratification to the derived formal nerve.}
Let $M$ be a coherent ordinary stratified module on $X/S$.  Continue locally with $X=\Spec A$ and $S=\Spec R$.  Let $Q_q$ be the completion of $A^{\otimes_R(q+1)}$ along the intrinsic diagonal.  Each projection $A\to Q_q$ is flat: tensoring the other flat $R$-factors is flat, and Noetherian completion is flat.  Therefore
\[
                    Q_q\otimes_A^{\mathbf L}M=Q_q\otimes_AM[0].
\]
These finite $Q_q$-modules are complete for the diagonal ideal.  The compatible ordinary stratification isomorphisms on the finite diagonal neighbourhoods give isomorphisms between these completed modules, because completion commutes with finite modules.  The identities and cocycle conditions continue to hold after completion.

These are derived stratification data as well.  Namely, derived pullback along a map between terms of the formal nerve is derived completed pullback; on a module of the displayed form it gives the corresponding $Q_r\otimes_A^{\mathbf L}M$, which is again the ordinary module because $Q_r$ is $A$-flat.  All mapping spaces between these modules are zero-truncated, so the ordinary identities and cocycles supply the coherent derived data.  This reasoning concerns the complete coherent modules; it does not assert that the finite diagonal neighbourhoods are flat over $A$.

Jiang's Theorem~2.1 and Remark~2.4 identify this intrinsic derived stratifying category; Proposition~6.2 identifies a formal neighbourhood with derived complete modules (equivalently supported modules); and Theorem~6.3 identifies the category with derived quasi-coherent infinitesimal crystals, compatibly with evaluation and restriction.  The hypotheses hold because $A$ is finitely presented and flat over the Noetherian ring $R$, hence has finite Tor-amplitude.  See \citet[pp.~16--17, 32--34; also Theorem~B.5, pp.~43--45]{Jiang}.  We obtain a derived crystal $K$ with $K_X=M[0]$ and with the prescribed completed stratification.

\emph{Zeroth cohomology gives an ordinary coherent crystal.}
For every finite-type test $U=\Spec C\hookrightarrow T=\Spec B$, its derived reduction satisfies
\[
 K_T\otimes_B^{\mathbf L}C=K_U\in D^{\leq0}(C).
\]
Indeed $K_U$ is the derived pullback of the ordinary module $M$ along $U\to X$.  For every $C$-module $Q$, associativity and right t-exactness give
\[
                   K_T\otimes_B^{\mathbf L}Q
                   =K_U\otimes_C^{\mathbf L}Q\in D^{\leq0}(B).
\]
Filter $B$ by the finitely many powers of the nilpotent ideal $I$.  Its successive quotients are $C$-modules.  The exact triangles of this finite filtration show that $K_T\in D^{\leq0}(B)$.  No bounded-below assumption has been used.

Set $G_T=H^0(K_T)$.  If $h:T\to T'$ is a morphism of tests, derived crystal compatibility and right t-exactness imply
\[
 H^0(K_T)=H^0(Lh^*K_{T'})=h^*H^0(K_{T'}).
\]
The last equality follows, for any complex in $D^{\leq0}$, by its truncation triangle: derived pullback sends the $D^{\leq-1}$ term into $D^{\leq-1}$.  Thus the $G_T$ form an ordinary quasi-coherent crystal.  Moreover
\[
          G_T/IG_T=H^0(K_U)=C\otimes_AM,
\]
which is a finite $C$-module.  Nilpotent Nakayama makes $G_T$ finite over $B$, and hence coherent on these Noetherian tests.  On $X$ its value is $M$, and on the finite intrinsic diagonal neighbourhoods its ordinary stratification is exactly the original one: zeroth cohomology of the derived reduction of the completed isomorphism is its ordinary reduction.  This proves essential surjectivity on finite-type tests.

\emph{Arbitrary tests and gluing.}
Work over $V=\Spec A$ and choose a presentation
\[ A=R[z_1,\ldots,z_r]/K_0. \]
Every affine test over $V$ maps to a Noetherian finite-type test
\[
        V\hookrightarrow\Spec\bigl(R[z_1,\ldots,z_r]/K_0^N\bigr)
\]
by lifting the $z_i$ and choosing $N$ large enough to annihilate the nilpotent ideal of the given test.  Any two choices admit a common refinement, obtained from their finite-type product by taking a sufficiently high graph neighbourhood of $V$.  Parallel compatible maps are equalized by their closed equalizer, which is still a Noetherian finite-type thickening of $V$.  Consequently pullback extends the ordinary crystal uniquely, on objects and morphisms, from finite-type tests to arbitrary classical tests.  The values on arbitrary tests are of finite presentation, because they are pullbacks of finite modules over Noetherian rings.

Full faithfulness, already proved locally, makes the local constructions for $M$ glue uniquely on overlaps of affine opens of $X$ and $S$; their cocycle relations follow from the same uniqueness.  The argument works on either infinitesimal site, with the domain of a big test mapping to the chosen affine open.  We have therefore obtained the claimed equivalence with ordinary tensor pullback in the full stated scope.
\end{proof}

The construction in this proof does not say that every derived test value has only degree-zero cohomology. Nonflat coefficients can produce negative cohomology. The upper bound is the precise fact needed for $H^0$ to supply an ordinary crystal.

\begin{proof}[Proof of Theorem~\ref{thm:flat}]
We state explicitly the relative derived comparison being used.  Zariski locally, $S$ and $X$ are affine Noetherian schemes and $X/S$ is separated, of finite type and of finite Tor-amplitude.  Jiang's Theorem~2.1 and Remark~2.4 identify the category of derived stratifications on the intrinsic formal diagonal nerve with modules over his derived differential-operator algebra.  His Theorem~6.3 identifies the latter category with derived quasi-coherent crystals on $(X/S)_{\rm inf}$, compatibly with evaluation on $X$; Appendix~B, Theorem~B.5, identifies the small and big versions \citep[pp.~16--17, 32--34 and 43--45]{Jiang}.

For clarity, the formal-diagonal interpretation uses Jiang's Proposition~6.2: on an affine closed immersion with ideal $J$, the category of derived sheaves on the formal completion is
\[
       \varprojlim_n D(B/J^n),
\]
with derived pullback transitions.  Equivalently it is the category of derived $J$-complete objects, or, by the local-cohomology/completion equivalence, the category denoted $\Gamma_J(\mathrm{QCoh}(\Spec B))$ in his Theorem~2.1.  The pullbacks in that theorem are the corresponding supported pullbacks.  Remark~2.4 gives precisely the resulting derived stratifying category.  We use this comparison only after checking that the ordinary data define objects and morphisms in these derived categories.  Flatness of $X/S$ ensures that its iterated products over $S$ here are the ordinary products, without extra derived structure.

Let $M$ be an ordinary flat stratified $\OO_X$-module.  On every finite diagonal neighbourhood, all its pullbacks are flat over that neighbourhood.  Their ordinary pullbacks therefore equal their derived pullbacks.  The given compatible isomorphisms and cocycle conditions define an object on the derived formal diagonal nerve.  Mapping spaces between the ordinary modules at each level have no positive homotopy groups, so these ordinary cocycle identities supply the required coherent identifications.  The cited comparison supplies a derived crystal $\mathcal K$ with $\mathcal K_X=M[0]$.

We now prove that its values are ordinary flat modules; this is a separate step from the cited comparison.  Let $B$ be any ring, let $I^N=0$, and let $K\in D(B)$ satisfy
\[
          K\otimes_B^{\mathbf L}B/I\simeq P[0],
          \qquad P\text{ flat over }B/I.
\]
For every $B/I$-module $Q$, tensor associativity gives
$K\otimes_B^{\mathbf L}Q\simeq(P\otimes_{B/I}Q)[0]$.
Filter an arbitrary $B$-module $Q'$ by the finite filtration $I^jQ'$.  Each successive quotient is a $B/I$-module, so the associated exact triangles show that $K\otimes_B^{\mathbf L}Q'$ is concentrated in degree zero.  Taking $Q'=B$ proves that $K$ is an ordinary module, and varying $Q'$ proves that it is flat.  This argument allows unbounded $K$ and imposes no Noetherian condition on $B$.

For a test $U\hookrightarrow T$, apply this argument to $K=\mathcal K_T$, whose derived reduction is the flat module pulled back from $M$ on $U$.  Thus $\mathcal K$ is an ordinary flat crystal.  Conversely an ordinary flat crystal has derived Cartesian transition maps, so defines a derived crystal.  The two constructions are inverse by the cited comparison, and the same reasoning on the zero-truncated mapping spaces identifies their ordinary morphisms.

These arguments initially use Jiang's finite-type classical tests.  They extend to arbitrary tests without an assumption about infinite completed tensor products.  Indeed, locally a test $\Spec C\hookrightarrow\Spec B$ mapping to an affine $\Spec A\subset X$, with $A=R[z_1,\ldots,z_r]/K$, maps to a finite-type test
\[
 \Spec A\hookrightarrow\Spec\bigl(R[z_1,\ldots,z_r]/K^N\bigr)
\]
by lifting the $z_i$, for sufficiently large $N$.  Two such choices have a common refinement: use the finite-type product of their ambient thickenings, then a sufficiently high infinitesimal neighbourhood of the diagonal copy of $\Spec A$.  Parallel compatible maps are equalized by their closed equalizer, which still contains that copy of $\Spec A$.  Pullback of the finite-type crystal is therefore independent of all choices and gives the unique extension, on objects and morphisms.  This also proves compatibility with affine localization and gluing.

Finally, if $M$ is finite locally free, nilpotent Nakayama shows that every test value is finitely generated.  At a local test ring, lift a basis of its reduction and use flatness to identify the kernel modulo the nilpotent ideal; a second nilpotent Nakayama argument makes that kernel zero.  The test value is consequently finite locally free.  The equivalence thus restricts as asserted.
\end{proof}

\begin{corollary}\label{cor:field}
For every field $k$ and every $k$-scheme $X$ locally of finite type, restriction gives an equivalence between coherent infinitesimal crystals and ordinary coherent stratified $\OO_X$-modules.  Smoothness, reducedness and $F$-purity are unnecessary.  Every coefficient module in either category is locally free of finite rank.
\end{corollary}
\begin{proof}
A coherent stratified module is locally free.  After the faithfully flat extension to an algebraic closure, fix a closed point $x$ and pull its stratification on each diagonal neighbourhood back along the second factor $x$.  The compatible isomorphisms identify the completed stalk with
\[
        \widehat{\OO}_{X,x}\otimes_k M_x/\mathfrak m_xM_x.
\]
Faithful flatness of the Noetherian local completion then makes the stalk free.  The nonfree locus of a finitely presented module is closed, so checking closed points on each affine finite-type open suffices.  Local freeness descends along the field extension.  This is also the argument behind Tonini--Zhang's Theorem~6.8(3).

For coherent infinitesimal crystals, Tonini--Zhang's Theorem~6.8(4), with its proof using Lemmas~6.10 and~6.15, shows that every test value is finite locally free \citep{ToniniZhang}.  The statement includes arbitrary ground fields.  Theorem~\ref{thm:flat} now applies, since a scheme locally of finite type over a field is flat over that field, and proves the asserted equivalence.
\end{proof}

\section{Relative Frobenius and arbitrary coefficients}

Affine space supplies a faithfully flat thickening at every Frobenius level. This replaces a local lift to a singular scheme by ordinary module descent, with no Noetherian requirement.

\begin{proof}[Proof of Theorem~\ref{thm:frobenius}]
We construct the equivalence locally and then identify its inverse on every infinitesimal test.  Locally on $X$ and $R$, choose a finitely presented closed immersion $X\hookrightarrow Y=\mathbf A_R^d$.  The iterated relative Frobenius $F_{Y/R}^n:Y\to Y^{(n)}$ is finite faithfully flat: in coordinates its module basis consists of the monomials with all exponents less than $p^n$.  More generally the same statement for smooth $Y/R$ is \cite[Lemma~50.19.5, Tag~0FW2]{StacksProject}.

Put $Y_n=Y\times_{Y^{(n)}}X^{(n)}$.  If $I=(f_1,\ldots,f_s)$ defines $X$ in $Y$, then $Y_n$ is defined by $(f_1^{p^n},\ldots,f_s^{p^n})$.  This remains true even though $F_R$ need not be flat: the ideal of a base-changed closed immersion is the image of the tensor product ideal, and relative Frobenius sends $f_i\otimes1$ to $f_i^{p^n}$.  Thus $X\hookrightarrow Y_n$ is nilpotent, and $Y_n\to X^{(n)}$ is finite faithfully flat.  Its double and triple fibre products are nilpotent thickenings of the diagonal copy of $X$.  Indeed, their defining difference ideals have finitely many generators killed by $p^n$th powers, in addition to the displayed nilpotent ideal.  A crystal therefore descends on this finite faithfully flat cover to a module $M_n$ on $X^{(n)}$.

The inclusion $Y_n\hookrightarrow Y_{n+1}$ and its maps on the double and triple products descend to $M_n\simeq F_{X^{(n)}/R}^*M_{n+1}$.  Conversely, for any nilpotent test $U\hookrightarrow T$ with ideal $J$, choose $n$ such that $J^{[p^n]}=0$.  In affine coordinates, the rule
\[
 (\OO_T/J)\otimes_{\OO_R,F_R^n}\OO_R
        \longrightarrow\OO_T,
       \qquad \overline b\otimes r\longmapsto b^{p^n}r
\]
is well defined and gives the relative Frobenius factor $\rho_n:T\to U^{(n)}$.  Set $\mathcal E_T=\rho_n^*(M_n|_{U^{(n)}})$.  The structural isomorphisms make this independent of $n$ and give the crystal transitions.

The inverse comparison is detected on
\[
                 Z=Y_n\times_{X^{(n)},\rho_n}T.
\]
It is finite faithfully flat over $T$ and is a nilpotent thickening of $U$.  For the latter assertion, modulo $J$ the extra difference ideal is finitely generated and killed elementwise by the $p^n$th power; adding $J$ preserves nilpotence.  This argument does not require $T$ to be Noetherian.  Ordinary faithfully flat module descent, including the double-product datum, proves that the constructions are inverse on all objects and morphisms.  Independence of the smooth embedding and gluing follow by the same canonical comparison.

Finally the relations $X\times_{X^{(n)}}X$ are cofinal among diagonal neighbourhoods: if $D$ is the diagonal ideal, their ideals are $D^{[p^n]}$, and, for $s$ generators of $D$, one has
\[D^{s(p^n-1)+1}\subseteq D^{[p^n]}\subseteq D^{p^n}.\]
The same bound applies to the triple diagonal ideals. Thus an ordinary stratification is exactly a compatible family of descent data at these levels.  Under the stated purity hypotheses all iterated relative Frobenius maps are pure.  Effective module descent \cite[Theorem~35.4.22, Tag~08WE]{StacksProject} supplies $M_n$ from the stratification at every level; full faithfulness supplies the transition identifications.  This proves the second assertion.
\end{proof}

\begin{corollary}\label{cor:node}
For every $\Fp$-algebra $R$, the flat finitely presented nodal family
\[
                    X=\Spec R[x,y]/(xy)\longrightarrow\Spec R
\]
satisfies the crystal--stratification equivalence for arbitrary quasi-coherent coefficients.
\end{corollary}
\begin{proof}
As an $R$-module its algebra is free on $1,x,x^2,\ldots,y,y^2,\ldots$, so the family is flat.  Every twist has the same displayed presentation.  The relative Frobenius $R[x_n,y_n]/(x_ny_n)\to R[x,y]/(xy)$ sends $x_n$ to $x^p$, $y_n$ to $y^p$, and fixes $R$.  It has an $R[x_n,y_n]/(x_ny_n)$-linear retraction: retain precisely the monomials whose exponent is divisible by $p$, divide that exponent by $p$, and keep their coefficients in $R$.  Linearity follows directly from $xy=0$.  Hence all relative Frobenius maps are split, in particular pure, and Theorem~\ref{thm:frobenius} applies.
\end{proof}

Over a perfect field, the pure-relative-Frobenius hypothesis includes every $F$-pure scheme of finite type: twists identify the required maps with iterates of absolute Frobenius. The nodal example proves a relative statement over every $\Fp$-algebra, including nonreduced and non-Noetherian bases.

\section{A crystal invisible on the underlying scheme}

The following formulas produce an infinitely compatible failure of faithfulness. A failure of descent at one Frobenius level alone would not suffice, because it might disappear at higher levels.

\begin{proof}[Proof of Theorem~\ref{thm:counterexample}]
Put
\[
 R_0=\Fp[t],\qquad I=\Fp[t,t^{-1}]/\Fp[t],\qquad
 R=R_0\ltimes I.
\]
Here the last expression is the square-zero extension with multiplication
$(a,u)(b,v)=(ab,av+bu)$.  We regard $t$ as $(t,0)$ and $I$ as the ideal $0\oplus I$ of $R$.  For $n\geq0$ set
\[
 s_n=1+p+\cdots+p^{n-1}=\frac{p^n-1}{p-1},\qquad
 b_n=(0,[t^{-s_n}])\in I,
\]
where $s_0=0$, so that $b_0=0$.  These elements satisfy
\begin{equation}\label{eq:recurrence}
 t^{p^n}b_{n+1}=b_n\quad(n\geq0),\qquad
 tb_1=0,\qquad b_1\ne0.
\end{equation}
Let
\[
 f=x^p-tx,\qquad A=R[x]/(f),\qquad X=\Spec A.
\]
The polynomial $f$ is monic.  Hence $A$ is a free $R$-module with basis $1,x,\ldots,x^{p-1}$; in particular $X/R$ is flat and finitely presented.  Moreover $f$ is a non-zero-divisor in $R[x]$, and every fibre is a zero-dimensional hypersurface over a field, hence a complete intersection and Cohen--Macaulay.  The example does not rely on a failure of the Cohen--Macaulay property of the fibres.

We construct a nonzero endomorphism $\delta$ of the structure crystal whose evaluation on $X$ is zero.  Let $U\hookrightarrow T$ be any infinitesimal test object, with nilpotent ideal $J$.  Locally on $T$, choose an integer $n$ such that $J^{[p^n]}=0$ and a lift $z\in\OO_T$ of the function $x|_U$.  Define
\begin{equation}\label{eq:phantom}
                     \delta_T=b_nz^{p^n}\in\OO_T.
\end{equation}
Changing $z$ by an element of $J$ does not change this function, because $T$ has characteristic $p$.  Also $z^p-tz\in J$, and therefore
\[
 b_{n+1}z^{p^{n+1}}-b_nz^{p^n}
       =b_{n+1}(z^p-tz)^{p^n}=0
\]
by \eqref{eq:recurrence}.  The definition is consequently independent of every sufficiently large $n$.  These two observations show that the local functions glue, including when only local nilpotence bounds have been chosen.

For a morphism of infinitesimal test objects, choose a common sufficiently large $n$ and pull back a local lift $z$.  Formula~\eqref{eq:phantom} is then preserved by pullback.  Thus the functions $\delta_T$ define a natural $\OO$-linear endomorphism of the structure crystal.  This construction works on all test objects, without a finiteness or quasi-coherence restriction on their nilpotent ideals.  It also works for the big infinitesimal site if that convention is adopted.

On the test object $X\hookrightarrow X$ one may take $n=0$, giving
$\delta_X=b_0x=0$.  On the other hand, take the square-zero test
\[
             X\hookrightarrow T=\Spec B,\qquad B=R[x]/(f^2).
\]
Here $J=(f)/(f^2)$ satisfies $J^2=0$, so $n=1$ is allowed for every $p$.  Formula~\eqref{eq:phantom} gives
\[
                  \delta_T=b_1x^p=b_1f\ne0.
\]
The equality uses $tb_1=0$.  Nonvanishing follows from the monicity of $f^2$: the classes of $1,x,\ldots,x^{2p-1}$ form an $R$-basis of $B$, and the coefficient of $x^p$ is $b_1\ne0$.  In fact this test $T$ is itself finite free over $R$.

Finally, $I^2=0$ implies $\delta_T^2=0$ for every $T$.  Hence $1+\delta$ is an automorphism of the rank-one structure crystal, with inverse $1-\delta$, and its restriction to the ordinary stratified module $\OO_X$ is the identity.  The restriction functor is not faithful, proving the assertion.
\end{proof}

\begin{corollary}[Failure over a regular Noetherian base with arbitrary quasi-coherent coefficients]\label{cor:noetherian-counterexample}
The restriction functor of Theorem~\ref{thm:counterexample} can fail to be faithful even when the base is the regular Noetherian ring $R=\Fp[t]$ and $X/R$ is finite free with complete-intersection fibres.  In this version both crystals are quasi-coherent; one is the rank-one structure crystal and the other has arbitrary, noncoherent coefficients.
\end{corollary}
\begin{proof}
Use $R=\Fp[t]$, the $R$-module $I=R[t^{-1}]/R$, and $A=R[x]/(x^p-tx)$.  We use no square-zero extension of the base.  Besides the structure crystal, consider the constant-base coefficient crystal
\[
                 \mathcal E_T=\OO_T\otimes_R I.
\]
Its crystal transition isomorphisms are ordinary tensor associativity, so no flatness assumption on $I$ is required.  Each value is a quasi-coherent module.

With $s_n$ and $b_n=[t^{-s_n}]\in I$ as before, the sections
\[
                    z^{p^n}\otimes b_n\in\mathcal E_T
\]
define a morphism $\OO_{\rm inf}\to\mathcal E$.  The lift-independence, compatibility in $n$, gluing, and naturality on all test objects are exactly the two identities proved in Theorem~\ref{thm:counterexample}, now interpreted in a module rather than a square-zero ideal.  Its value on $X$ is zero, whereas its value on $B=R[x]/(x^p-tx)^2$ is
\[
                         x^p\otimes b_1\ne0.
\]
Indeed, $B\otimes_R I$ is the direct sum of $2p$ copies of $I$ in the monic basis, and its $x^p$ coefficient is $b_1\ne0$.  This proves nonfaithfulness.  Equivalently, the upper-triangular automorphism of $\OO_{\rm inf}\oplus\mathcal E$ with this off-diagonal entry is a nonidentity automorphism invisible to restriction.

The module $I$ is not finitely generated over $R$.  Since $A$ is finite faithfully flat and free over $R$, the module $A\otimes_R I$ is not finitely generated over $A$: otherwise it would be finitely generated over $R$, as would its direct summand $I$.  Thus this example does not fall under the coherent-coefficient hypothesis.
\end{proof}

For the smallest example take $p=2$. Then $s_0=0$, $s_1=1$, $s_2=3$ and $s_3=7$. On $R[x]/(x^2-tx)^2$, the invisible endomorphism has value $[t^{-1}]x^2$. Its coefficient is nonzero in the monic basis, although its value on $X$ is zero. The program \emph{Crystal principal parts}, available through the workbench, checks the first recurrence levels and these finite witnesses by exact Laurent-exponent arithmetic. The identities $t^{2^n}[t^{-s_{n+1}}]=[t^{-s_n}]$ explain why this finite witness extends compatibly to all thicker tests.

\section{Cuspidal normalizations}

The cusp gives a useful boundary test: its Frobenius is not pure, yet arbitrary coefficient stratifications still determine crystals. We prove this by ordinary faithfully flat descent, with no flatness requirement on a coefficient module.

\begin{lemma}\label{lem:homeomorphism}
Over a perfect field of characteristic $p$, pullback along a finite universal homeomorphism between reduced schemes locally of finite type induces an equivalence on the categories of $F$-divided quasi-coherent modules, with arbitrary quasi-coherent coefficients.
\end{lemma}
\begin{proof}
Write $f:Y\to X$ for the morphism. The assertion is local on the target. For an affine such morphism write $A\subset B$; injectivity follows because the schemes are reduced and the map is surjective. We first prove that $B^{p^N}\subset A$ for some $N$. This power property is also the content of \cite[Lemma~29.47.9, Tag~0CNF]{StacksProject} in this setting.

Write $Q(A)$ for the total quotient ring of $A$. The total quotient rings of the reduced algebras are products of their component fields. A finite universal homeomorphism pairs the components and gives purely inseparable extensions of their fields. Choose finitely many algebra generators for $B/A$ and an exponent $r$ for which their $p^r$-powers lie in the total quotient ring of $A$. Then
\[
 C=A[B^{p^r}]\subset Q(A)
\]
is finite and birational over $A$ and still a universal homeomorphism. Its conductor $J=\{a\in A:aC\subset A\}$ is an ideal of both rings and contains a nonzerodivisor, obtained by clearing denominators of finitely many module generators. Consequently $A/J\hookrightarrow C/J$ has smaller dimension on every component. Induction on dimension gives an exponent $s$ for which every $p^s$-power in $C/J$ belongs to $A/J$. Since $J\subset A$, every $p^s$-power in $C$ belongs to $A$, and $B^{p^{r+s}}\subset A$. The induction also permits nonreduced quotient rings: first reduce them, then increase the exponent to kill their nilpotent nilradicals. In dimension zero it is the elementary power property for finite purely inseparable field extensions, followed by the same nilpotent reduction.

The power map gives a relative Frobenius factor
\[
 g:X\longrightarrow Y^{(N)},\qquad
 gf=F_{Y/k}^N,\quad f^{(N)}g=F_{X/k}^N.
\]
On rings it sends $b\otimes c$ to $b^{p^N}c$. For an $F$-divided system $N_\bullet$ on $Y$, set
\[
 Q(N)_r=(g^{(r)})^*N_{r+N}.
\]
The transition maps come from those of $N_\bullet$ and the commutation with relative Frobenius. The two factorization identities show
\[
 (f^{(r)})^*Q(N)_r\simeq N_r,\qquad
 Q(f^*M)_r\simeq M_r.
\]
These isomorphisms respect all transitions and morphisms. They give inverse equivalences without any condition on the modules. The fully faithful local inverses glue, proving the assertion.
\end{proof}

\begin{lemma}\label{lem:flat-homeomorphism}
Let $h:V\to W$ be a finite faithfully flat universal homeomorphism of separated finite-type schemes over a field. Ordinary pullback is an equivalence between their categories of stratified quasi-coherent modules, with arbitrary quasi-coherent coefficients.
\end{lemma}
\begin{proof}
The diagonal $V\hookrightarrow V\times_WV$ is a nilpotent closed immersion: the morphism is separated and universally injective, and the rings are Noetherian. A stratification on a quasi-coherent module $M$ on $V$ therefore restricts to a descent datum for $h$. Its cocycle is the ordinary triple identity. Faithfully flat descent gives a module $N$ on $W$ and an identification $h^*N\simeq M$.

Let $D$ be a finite neighbourhood of the diagonal in $W\times W$ and put
\[
 D'=D\times_{W\times W}(V\times V).
\]
The map $D'\to D$ is finite faithfully flat. The diagonal $V\hookrightarrow D'$ is nilpotent, since it is a closed immersion onto the underlying space and the schemes are Noetherian. Hence $D'$ factors through a finite diagonal neighbourhood in $V\times V$, where the given stratification compares the two pullbacks of $M$.

This comparison is compatible with descent on $D'\times_DD'$. That scheme describes four points of $V$, with the first and third images in $W$ equal, and likewise the second and fourth. It is supported on the small diagonal. The ordinary triple cocycle identities, pulled back to this quadruple neighbourhood, say that its comparison commutes with the two descent identifications. The comparison therefore descends to $D$. Identity and triple compatibility for the descended comparisons can be checked after the corresponding faithfully flat covers by products of $V$. They hold there by the given stratification, so they hold on $W$.

Morphisms descend in the same way. Starting with a stratified module on $W$, the construction recovers its canonical descent datum after pullback. These constructions are inverse. Only $h$ has been required to be flat; the coefficient modules were arbitrary.
\end{proof}

\begin{proposition}\label{cor:cusps}
Let $k$ be a perfect field of characteristic $p$, and let $X$ be an integral Cohen--Macaulay $k$-variety admitting a finite universal homeomorphism $f:Y\to X$ with $Y$ smooth.  Restriction is an equivalence between quasi-coherent crystals on $(X/k)_{\rm inf}$ and ordinary stratified quasi-coherent $\OO_X$-modules, with no finiteness or flatness requirement on the coefficients.

In particular this applies to the cusp $X=\Spec k[t^2,t^3]$, although this cusp is not $F$-pure for any prime $p$.
\end{proposition}
\begin{proof}
On a finite affine cover the preceding power argument gives a common exponent $N$ and a map $g:X\to Y^{(N)}$ with $gf=F_{Y/k}^N$ and $f^{(N)}g=F_{X/k}^N$. These maps glue because their ring maps are the same canonical power maps on overlaps. The map $g$ is finite: $X$ is finite over its Frobenius twist, since $k$ is perfect and $X$ is of finite type, and the image of the target algebra of $g$ contains the image of this Frobenius algebra. A finite set of generators over the smaller algebra therefore also generates over the larger one. The factorization shows that $g$ is a universal homeomorphism.

For every point $x$ of $X$, the local ring of $Y^{(N)}$ at $g(x)$ is regular, the local ring of $X$ at $x$ is Cohen--Macaulay, their dimensions agree, and the fibre has dimension zero. Miracle flatness \cite[Tag~00R4]{StacksProject} implies that $g$ is flat. It is surjective, so it is finite faithfully flat. This is the precise use of Cohen--Macaulayness; the original map $f$ need not be flat.

By Lemma~\ref{lem:flat-homeomorphism}, $g^*$ is an equivalence on arbitrary quasi-coherent stratifications. It is also an equivalence on quasi-coherent infinitesimal crystals by Theorem~\ref{thm:frobenius} and Lemma~\ref{lem:homeomorphism}. On the smooth scheme $Y^{(N)}$, relative Frobenius is finite locally free and surjective, as seen from an \'etale coordinate chart and the monomial basis on affine space. Its purity and Theorem~\ref{thm:frobenius} give the crystal--stratification comparison there. Restriction commutes with $g^*$, so the comparison on $X$ follows from these three equivalences.

For comparison, Ben-Zvi--Nevins \cite[Proposition~2.11, Theorems~3.11 and~4.3, Remark~4.5]{BenZviNevins} prove the ordinary-pullback equivalence along cuspidal quotients under these hypotheses. Their Remark~2.7 uses tensor products in the pro-category, corresponding to compatible ordinary stratifications on all finite diagonal neighbourhoods. The direct descent proof here gives the coefficient scope needed for the stated comparison.

For the stated cusp, $k[t]$ is its finite normalization, and the map is bijective on geometric points, hence a universal homeomorphism.  The curve is a hypersurface and therefore Cohen--Macaulay.  To verify nonpurity directly, put $A=k[t^2,t^3]$ and $I=(t^2)$.  The class of $t^3$ in $A/I$ is nonzero, whereas its image under the Frobenius test map is zero:
\[
       t^3\notin(t^2),\qquad
       (t^3)^p=t^{3p}\in(t^{2p})=I^{[p]},
\]
because $t^p\in A$ for every $p\geq2$.  Thus the Frobenius map is not pure, and this positive case is not a consequence of Theorem~\ref{thm:frobenius}.
\end{proof}

\section{The scope of the comparison}

The Noetherian coherent comparison is complete under the hypotheses of Theorem~\ref{thm:coherent}. For arbitrary coefficients, a useful remaining problem is to characterize flat morphisms of finite presentation for which ordinary restriction is an equivalence. Theorem~\ref{thm:frobenius} gives one sufficient condition in characteristic $p$, and the cusp shows that this condition is not necessary. The counterexample over $\Fp[t]$ prevents replacing such a criterion by Noetherianity of the base alone. Its obstruction is an infinite Frobenius-compatible family of sections invisible at level zero.

A second question is to identify weaker finiteness conditions on a non-Noetherian base that still give the coherent comparison. The proof here uses flatness of adic completion and exact completion of finite modules over Noetherian rings. The square-zero principal-parts base demonstrates a concrete failure when those inputs are absent. No extension to every coherent or weakly proregular non-Noetherian ring is asserted.

\section*{Note on method and reuse}

The mathematics and text were produced by GPT-6 Astra Pro, the model displayed by the interface, in ChatGPT Pro mode. The written arguments and the indicated primary-source statements received targeted checks by the same model. Separate statement-only derivations by the same model reconstructed all six main statements from a proof-unseen packet and were frozen before comparison. Those proofs were then checked in full; they include an ordinary completion-and-descent proof of the Noetherian assertions and the finite-flat Frobenius-factor proof used above. Effort settings are not exposed by the interface. No human has checked this paper. These statements distinguish the checks performed from independent verification by another model family. The text and original computations are dedicated to the public domain under CC0~1.0. The verification records and programs are maintained through \url{\workbenchurl}; the permanent record address is \url{\recordurl}.

\begin{thebibliography}{99}
\bibitem[Ben-Zvi and Nevins(2004)]{BenZviNevins}
Ben-Zvi, D., \& Nevins, T. (2004). Cusps and $\mathcal D$-modules. \emph{Journal of the American Mathematical Society, 17}(1), 155--179. \url{https://doi.org/10.1090/S0894-0347-03-00439-9}. Statement numbers and preprint pages refer to \url{https://arxiv.org/abs/math/0212094v3}.
\bibitem[Grothendieck(n.d.)]{GrothendieckNotes}
Grothendieck, A. (n.d.). \emph{Lettre Tate (mai 66) (Cristaux)}. Archives Grothendieck, Universit\'e de Montpellier, cote~6, \S1.8, archive sheet~6 (PDF image~7). \url{https://grothendieck.umontpellier.fr/6.pdf}.
\bibitem[Grothendieck(1968)]{Grothendieck1968}
Grothendieck, A. (1968). Crystals and the de Rham cohomology of schemes. In \emph{Dix expos\'es sur la cohomologie des sch\'emas} (pp.~306--358). North-Holland and Masson. Notes by J. Coates and O. Jussila. \url{https://download.uni-mainz.de/mathematik/Algebraische%20Geometrie/Lehre/WS23.Padische.1966.Grothendieck.CRCSscan.pdf}.
\bibitem[Hua(2026)]{Hua}
Hua, M. (2026). \emph{\'Enonc\'es des notes de travail de Grothendieck d\'emontr\'es en Lean4, et questions ouvertes} (Version~2). Hexagon, 2611.00125. \url{https://hexagonmath.org/2611.00125v2}.
\bibitem[Jiang(2023)]{Jiang}
Jiang, A. (2023). \emph{The derived ring of differential operators}. arXiv:2303.16083v1. \url{https://arxiv.org/abs/2303.16083}.
\bibitem[Mesablishvili(2000)]{Mesablishvili}
Mesablishvili, B. (2000). Pure morphisms of commutative rings are effective descent morphisms for modules---a new proof. \emph{Theory and Applications of Categories, 7}(3), 38--42. \url{https://ftp.gwdg.de/pub/misc/EMIS/journals/TAC/volumes/7/n3/n3.pdf}.
\bibitem[The Stacks Project Authors(2026)]{StacksProject}
The Stacks Project Authors. (2026). \emph{The Stacks Project}. Tags 00IP, 00IQ, 00HR, 00MA, 00MB, 00R4, 08WE, 0CNF, and 0FW2. \url{https://stacks.math.columbia.edu}.
\bibitem[Tonini and Zhang(2019)]{ToniniZhang}
Tonini, F., \& Zhang, L. (2019). Algebraic and Nori fundamental gerbes. \emph{Journal of the Institute of Mathematics of Jussieu, 18}(4), 855--897. \url{https://doi.org/10.1017/S147474801700024X}. Preprint locators refer to \url{https://arxiv.org/abs/1605.01332v2}.
\end{thebibliography}
\end{document}
